\documentclass[12pt,a4paper]{article}
\usepackage[utf8]{inputenc}
\usepackage[spanish]{babel}
\usepackage{geometry}
\usepackage{fancyhdr}
\usepackage{enumitem}
\usepackage{xcolor}
\usepackage{listings}
\usepackage{graphicx}
\usepackage{hyperref}
\usepackage{tabularx}
\usepackage{array}

\geometry{margin=2cm}
\pagestyle{fancy}
\fancyhf{}
\lhead{Curso de Programación en Python}
\rhead{Semana 4}
\cfoot{\thepage}

% Configuración para código Python
\lstset{
    language=Python,
    basicstyle=\ttfamily\small,
    keywordstyle=\color{blue},
    stringstyle=\color{red},
    commentstyle=\color{green!60!black},
    showstringspaces=false,
    breaklines=true,
    frame=single,
    numbers=left,
    numberstyle=\tiny\color{gray}
}

\title{\textbf{Tarea Semana 4: Simulador de Lotería Nacional}\\
\large Librerías y Módulos en Python}
\author{Curso de Programación en Python}
\date{Tiempo de entrega: 1 semana}

\begin{document}

\maketitle

\section{Introducción}

¡Bienvenido a la Semana 4! Has aprendido los fundamentos de Python: variables, condicionales, bucles y manejo de errores. Ahora es momento de ampliar tu poder de programación usando \textbf{librerías y módulos}.

Las librerías son colecciones de código pre-escrito que puedes usar en tus programas. ¿Por qué reinventar la rueda cuando alguien ya creó una solución perfecta? Python tiene una vasta biblioteca estándar y miles de paquetes externos disponibles.

En esta tarea, crearás un simulador completo de la Lotería Nacional de Costa Rica, usando múltiples librerías para generar números aleatorios, calcular estadísticas, trabajar con fechas, interactuar con el sistema operativo, y más.

\section{Objetivos de Aprendizaje}

Al completar esta tarea, serás capaz de:

\begin{itemize}[leftmargin=*]
    \item Importar módulos con \texttt{import}
    \item Usar importación selectiva con \texttt{from ... import ...}
    \item Trabajar con módulos de la biblioteca estándar de Python
    \item Usar \texttt{random} para generación de números aleatorios
    \item Usar \texttt{statistics} para cálculos estadísticos
    \item Usar \texttt{datetime} para manejo de fechas y horas
    \item Usar \texttt{os} y \texttt{sys} para interacción con el sistema
    \item Instalar paquetes externos con \texttt{pip}
    \item Crear tus propios módulos personalizados
    \item Organizar código en múltiples archivos
    \item Entender la diferencia entre módulos y paquetes
\end{itemize}

\section{Enunciado del Proyecto}

\subsection{Descripción General}

Vas a crear un programa llamado \textbf{``Simulador de Lotería Nacional de Costa Rica''}. Este programa simula el sorteo de la Lotería Nacional, permite a usuarios comprar billetes virtuales, realizar sorteos, calcular premios, mantener estadísticas, y guardar/cargar datos.

\subsection{Contexto de la Lotería Nacional}

La Lotería Nacional de Costa Rica funciona así:
\begin{itemize}
    \item Los billetes tienen números de 00000 a 99999 (5 dígitos)
    \item Cada billete tiene una serie (A, B, C, etc.)
    \item El premio mayor es para el número exacto
    \item Hay premios para aproximaciones (un número antes/después)
    \item Hay premios para terminaciones (últimos 2, 3, 4 dígitos)
    \item Los sorteos se realizan en fechas específicas
\end{itemize}

\section{Requisitos Funcionales}

\subsection{Parte 1: Módulos a Utilizar}

Tu programa DEBE usar las siguientes librerías:

\subsubsection{1. Módulo random}

\begin{lstlisting}[caption=Uso de random para sorteos]
import random

# Generar numero ganador (00000-99999)
numero_ganador = random.randint(0, 99999)

# Generar serie ganadora (A-Z)
series = ['A', 'B', 'C', 'D', 'E']
serie_ganadora = random.choice(series)

# Simular venta de billetes aleatorios
billetes_vendidos = random.sample(range(100000), 1000)

# Generar multiples ganadores para aproximaciones
aproximaciones = [numero_ganador - 1, numero_ganador + 1]
\end{lstlisting}

\subsubsection{2. Módulo statistics}

\begin{lstlisting}[caption=Uso de statistics para analisis]
import statistics

# Lista de precios de billetes vendidos
precios = [10000, 10000, 10000, 5000, 5000]

# Calcular estadisticas de ventas
promedio_precio = statistics.mean(precios)
mediana_precio = statistics.median(precios)
moda_precio = statistics.mode(precios)

# Calcular desviacion estandar de numeros jugados
numeros_jugados = [12345, 54321, 98765, 11111]
desv_std = statistics.stdev(numeros_jugados)
\end{lstlisting}

\subsubsection{3. Módulo datetime}

\begin{lstlisting}[caption=Uso de datetime para fechas]
from datetime import datetime, timedelta

# Obtener fecha y hora actual
fecha_actual = datetime.now()
print(f"Sorteo realizado: {fecha_actual.strftime('%d/%m/%Y %H:%M:%S')}")

# Calcular proxima fecha de sorteo (cada domingo)
dias_hasta_domingo = (6 - fecha_actual.weekday()) % 7
proxima_fecha = fecha_actual + timedelta(days=dias_hasta_domingo)

# Registrar tiempo de sorteo
inicio = datetime.now()
# ... realizar sorteo ...
fin = datetime.now()
duracion = (fin - inicio).total_seconds()
print(f"Sorteo completo en {duracion} segundos")
\end{lstlisting}

\subsubsection{4. Módulo os}

\begin{lstlisting}[caption=Uso de os para sistema de archivos]
import os

# Verificar si existe carpeta de datos
if not os.path.exists('datos_loteria'):
    os.makedirs('datos_loteria')
    print("Carpeta de datos creada")

# Listar archivos de sorteos anteriores
archivos = os.listdir('datos_loteria')
print(f"Sorteos guardados: {len(archivos)}")

# Obtener informacion del sistema
print(f"Sistema operativo: {os.name}")
print(f"Directorio actual: {os.getcwd()}")
\end{lstlisting}

\subsubsection{5. Módulo sys}

\begin{lstlisting}[caption=Uso de sys para informacion del sistema]
import sys

# Mostrar version de Python
print(f"Python version: {sys.version}")

# Verificar argumentos de linea de comandos
if len(sys.argv) > 1:
    modo = sys.argv[1]
    print(f"Modo: {modo}")

# Salir del programa con codigo de estado
if error_critico:
    print("Error critico detectado")
    sys.exit(1)
\end{lstlisting}

\subsubsection{6. Librería externa: colorama (opcional)}

\begin{lstlisting}[caption=Instalacion y uso de colorama para colores]
# Instalar con: pip install colorama

from colorama import Fore, Back, Style, init

# Inicializar colorama
init(autoreset=True)

# Imprimir mensajes con colores
print(Fore.GREEN + "FELICIDADES! HAS GANADO!")
print(Fore.RED + "Lo sentimos, no ganaste esta vez")
print(Fore.YELLOW + "Numero ganador: " + Fore.CYAN + "12345")
print(Back.WHITE + Fore.BLACK + "Premio Mayor")
\end{lstlisting}

\subsection{Parte 2: Crear Módulo Propio}

Debes crear al menos UN módulo personalizado llamado \texttt{utilidades\_loteria.py}:

\begin{lstlisting}[caption=Archivo utilidades\_loteria.py]
"""
Modulo de utilidades para el simulador de loteria
Contiene funciones auxiliares reutilizables
"""

def formatear_numero(numero):
    """Formatea un numero de loteria con ceros a la izquierda"""
    return str(numero).zfill(5)

def validar_numero_loteria(numero):
    """Valida que un numero este en rango valido (0-99999)"""
    try:
        num = int(numero)
        return 0 <= num <= 99999
    except ValueError:
        return False

def calcular_premio(numero_jugado, numero_ganador, tipo="exacto"):
    """Calcula el premio segun tipo de acierto"""
    premios = {
        "exacto": 200000000,  # 200 millones
        "aproximacion": 50000000,  # 50 millones
        "3_cifras": 500000,  # 500 mil
        "2_cifras": 100000   # 100 mil
    }
    return premios.get(tipo, 0)

def obtener_terminaciones(numero):
    """Obtiene las terminaciones de 2, 3 y 4 cifras"""
    num_str = str(numero).zfill(5)
    return {
        "2_cifras": num_str[-2:],
        "3_cifras": num_str[-3:],
        "4_cifras": num_str[-4:]
    }
\end{lstlisting}

Luego importar y usar en tu programa principal:

\begin{lstlisting}[caption=Importar modulo propio en main.py]
# Importar modulo completo
import utilidades_loteria

numero = utilidades_loteria.formatear_numero(1234)

# O importar funciones especificas
from utilidades_loteria import validar_numero_loteria, calcular_premio

if validar_numero_loteria(numero_jugado):
    premio = calcular_premio(numero_jugado, ganador, "exacto")
\end{lstlisting}

\subsection{Parte 3: Funcionalidades del Programa}

\subsubsection{Menú Principal}

\begin{enumerate}[leftmargin=*]
    \item Comprar billetes de lotería
    \item Realizar sorteo
    \item Ver mis billetes
    \item Ver resultados del último sorteo
    \item Ver estadísticas generales
    \item Verificar si gané
    \item Información del sistema
    \item Salir
\end{enumerate}

\subsubsection{Comprar Billetes}

\begin{itemize}
    \item Generar número aleatorio o elegir número específico
    \item Elegir serie (A-E) o aleatoria
    \item Cada billete cuesta ₡10,000
    \item Permitir comprar múltiples billetes
    \item Usar \texttt{random.randint()} para números aleatorios
    \item Guardar billetes del usuario en una lista
    \item Mostrar resumen de compra con fecha (\texttt{datetime})
\end{itemize}

\subsubsection{Realizar Sorteo}

\begin{itemize}
    \item Generar número ganador (00000-99999) con \texttt{random}
    \item Generar serie ganadora con \texttt{random.choice()}
    \item Generar aproximaciones (número -1 y +1)
    \item Registrar fecha y hora del sorteo con \texttt{datetime}
    \item Calcular todos los premios:
    \begin{itemize}
        \item Premio mayor (número exacto + serie)
        \item Aproximaciones
        \item 4 cifras (últimos 4 dígitos)
        \item 3 cifras (últimos 3 dígitos)
        \item 2 cifras (últimos 2 dígitos)
    \end{itemize}
    \item Guardar resultados del sorteo
\end{itemize}

\subsubsection{Verificar Premios}

\begin{itemize}
    \item Comparar billetes del usuario con resultados
    \item Usar módulo propio para calcular premios
    \item Mostrar si ganó y cuánto
    \item Actualizar estadísticas de ganancias
\end{itemize}

\subsubsection{Estadísticas}

Usar \texttt{statistics} para calcular:
\begin{itemize}
    \item Total de billetes comprados
    \item Total invertido
    \item Total ganado
    \item Promedio de números jugados
    \item Número más jugado (moda)
    \item Serie más jugada
    \item Rentabilidad (ganado/invertido)
    \item Historial de sorteos realizados
\end{itemize}

\subsubsection{Información del Sistema}

Usar \texttt{os} y \texttt{sys} para mostrar:
\begin{itemize}
    \item Versión de Python
    \item Sistema operativo
    \item Directorio actual
    \item Archivos guardados
    \item Espacio usado por datos
\end{itemize}

\newpage

\subsection{Ejemplo de Ejecución}

\begin{lstlisting}[caption=Flujo completo del programa]
========================================
  SIMULADOR LOTERIA NACIONAL CR
========================================
Python 3.11.5 | Sistema: Windows
Fecha: 24/10/2025 14:30:15

--- MENU PRINCIPAL ---
1. Comprar billetes
2. Realizar sorteo
3. Ver mis billetes
4. Resultados ultimo sorteo
5. Estadisticas
6. Verificar premios
7. Info del sistema
8. Salir

Opcion: 1

--- COMPRAR BILLETES ---
1. Numero aleatorio
2. Elegir numero especifico

Opcion: 1

Generando numero aleatorio...
Numero: 45678
Serie: C
Precio: ₡10,000

Confirmar compra? (s/n): s

Billete comprado exitosamente!
Fecha de compra: 24/10/2025 14:31:00

Comprar otro? (s/n): s

Opcion: 2
Ingresa el numero (00000-99999): 12345
Serie (A-E) o random (r): A
Precio: ₡10,000

Billete comprado exitosamente!

Comprar otro? (s/n): n

Total invertido: ₡20,000
Billetes comprados: 2

--- MENU PRINCIPAL ---
Opcion: 3

--- MIS BILLETES ---
1. Numero: 45678 | Serie: C | Comprado: 24/10/2025 14:31
2. Numero: 12345 | Serie: A | Comprado: 24/10/2025 14:32

Total: 2 billetes

--- MENU PRINCIPAL ---
Opcion: 2

--- REALIZANDO SORTEO ---
Fecha del sorteo: 24/10/2025 14:35:00

Generando numero ganador...
Mezclando bolas...
Seleccionando serie...

NUMERO GANADOR: 12346
SERIE GANADORA: A

Aproximaciones:
- 12345 (aproximacion anterior)
- 12347 (aproximacion posterior)

Premios por terminaciones:
- 4 cifras (2346): ₡2,000,000
- 3 cifras (346): ₡500,000
- 2 cifras (46): ₡100,000

Sorteo completado en 0.045 segundos

--- MENU PRINCIPAL ---
Opcion: 6

--- VERIFICAR PREMIOS ---
Revisando tus 2 billetes...

Billete #1: 45678 - Serie C
  Terminacion 2 cifras: NO
  Terminacion 3 cifras: NO
  Sin premio

Billete #2: 12345 - Serie A
  APROXIMACION ANTERIOR!
  Premio: ₡50,000,000

FELICIDADES! HAS GANADO ₡50,000,000!

Tu inversion: ₡20,000
Tus ganancias: ₡50,000,000
Rentabilidad: 250,000%

--- MENU PRINCIPAL ---
Opcion: 5

--- ESTADISTICAS GENERALES ---

Datos de Usuario:
- Total billetes comprados: 2
- Total invertido: ₡20,000
- Total ganado: ₡50,000,000
- Rentabilidad: 250,000%

Analisis de Numeros Jugados:
- Numeros: [45678, 12345]
- Promedio: 29,011.5
- Numero mas jugado: N/A (todos diferentes)

Series Jugadas:
- Serie A: 1 vez
- Serie C: 1 vez

Historial de Sorteos:
- Sorteos realizados: 1
- Ultimo sorteo: 24/10/2025 14:35
- Numero ganador: 12346 - Serie A

--- MENU PRINCIPAL ---
Opcion: 7

--- INFORMACION DEL SISTEMA ---

Python:
- Version: 3.11.5
- Path: C:\Python311\python.exe

Sistema Operativo:
- OS: Windows (nt)
- Directorio actual: C:\Users\user\loteria

Archivos:
- Sorteos guardados: 1
- Billetes guardados: 2
- Carpeta datos: datos_loteria/

Modulos cargados:
- random
- statistics
- datetime
- os
- sys
- utilidades_loteria

--- MENU PRINCIPAL ---
Opcion: 8

Guardando datos...
Gracias por usar el Simulador de Loteria!

Hasta luego!
\end{lstlisting}

\newpage

\section{Requisitos Técnicos Detallados}

\subsection{Importaciones Requeridas}

Tu archivo principal debe incluir:

\begin{lstlisting}[caption=Importaciones del programa principal]
# Modulos de biblioteca estandar
import random
import statistics
import os
import sys
from datetime import datetime, timedelta

# Modulo propio
import utilidades_loteria
# O importacion selectiva:
# from utilidades_loteria import formatear_numero, calcular_premio

# Libreria externa (opcional pero recomendado)
try:
    from colorama import Fore, Style, init
    init(autoreset=True)
    COLORAMA_DISPONIBLE = True
except ImportError:
    print("Colorama no instalado. Ejecutar: pip install colorama")
    COLORAMA_DISPONIBLE = False
\end{lstlisting}

\subsection{Estructura de Archivos}

Tu proyecto debe tener esta estructura:

\begin{lstlisting}[caption=Estructura del proyecto]
loteria/
    main.py                    # Programa principal
    utilidades_loteria.py      # Modulo personalizado
    datos_loteria/             # Carpeta para datos (se crea automaticamente)
        sorteos.txt
        billetes.txt
        estadisticas.txt
    requirements.txt           # Lista de dependencias
    README.txt                 # Instrucciones
\end{lstlisting}

\subsection{Archivo requirements.txt}

Crear archivo con dependencias:

\begin{lstlisting}[caption=requirements.txt]
colorama>=0.4.6
\end{lstlisting}

Instrucciones para instalar:
\begin{lstlisting}[language=bash]
pip install -r requirements.txt
\end{lstlisting}

\subsection{Uso Correcto de Módulos}

\subsubsection{Random}

Debes usar al menos 3 funciones diferentes:

\begin{lstlisting}
import random

# randint: numero entero aleatorio en rango
numero = random.randint(0, 99999)

# choice: elegir elemento aleatorio de lista
serie = random.choice(['A', 'B', 'C', 'D', 'E'])

# sample: multiples elementos sin repeticion
numeros_vendidos = random.sample(range(100000), 50)

# shuffle: mezclar lista (in-place)
lista = [1, 2, 3, 4, 5]
random.shuffle(lista)

# random: float entre 0 y 1
probabilidad = random.random()
if probabilidad < 0.01:  # 1% de probabilidad
    print("Bonus especial!")
\end{lstlisting}

\subsubsection{Statistics}

Debes usar al menos 3 funciones:

\begin{lstlisting}
import statistics

numeros = [12345, 54321, 11111, 99999, 12345]

# mean: promedio
promedio = statistics.mean(numeros)

# median: mediana
mediana = statistics.median(numeros)

# mode: moda (mas frecuente)
moda = statistics.mode(numeros)

# stdev: desviacion estandar
desviacion = statistics.stdev(numeros)

# variance: varianza
varianza = statistics.variance(numeros)
\end{lstlisting}

\subsubsection{Datetime}

Debes usar funcionalidades de fecha/hora:

\begin{lstlisting}
from datetime import datetime, timedelta

# Obtener fecha/hora actual
ahora = datetime.now()

# Formatear fecha
fecha_str = ahora.strftime("%d/%m/%Y %H:%M:%S")

# Parsear fecha desde string
fecha = datetime.strptime("24/10/2025", "%d/%m/%Y")

# Calcular diferencias de tiempo
proximo_sorteo = ahora + timedelta(days=7)
diferencia = proximo_sorteo - ahora
dias_restantes = diferencia.days

# Obtener componentes
anio = ahora.year
mes = ahora.month
dia = ahora.day
hora = ahora.hour
\end{lstlisting}

\subsection{Validaciones con Try-Except}

Todas las importaciones y operaciones deben manejarse:

\begin{lstlisting}
# Manejar modulos opcionales
try:
    from colorama import Fore
    TIENE_COLOR = True
except ImportError:
    TIENE_COLOR = False
    print("Instalando sin colores...")

# Manejar operaciones de archivo
try:
    with open('datos.txt', 'r') as f:
        datos = f.read()
except FileNotFoundError:
    print("Archivo no encontrado, creando nuevo...")
    datos = ""
except PermissionError:
    print("Sin permisos para leer archivo")
\end{lstlisting}

\section{Consejos y Recomendaciones}

\begin{itemize}[leftmargin=*]
    \item \textbf{Lee la documentación:} Cada módulo tiene documentación oficial. Úsala.
    
    \item \textbf{Importa solo lo necesario:} No hagas \texttt{from modulo import *}
    
    \item \textbf{Organiza las importaciones:}
    \begin{enumerate}
        \item Biblioteca estándar primero
        \item Librerías externas después
        \item Módulos propios al final
    \end{enumerate}
    
    \item \textbf{Usa alias cuando sea útil:} \texttt{import statistics as stats}
    
    \item \textbf{Comenta módulos propios:} Usa docstrings en funciones
    
    \item \textbf{Prueba instalación de paquetes:} Verifica que \texttt{pip} funciona
    
    \item \textbf{Maneja ImportError:} Los módulos externos pueden no estar instalados
    
    \item \textbf{No reinventes la rueda:} Si existe un módulo que hace lo que necesitas, úsalo
\end{itemize}

\section{Entregables}

\begin{enumerate}[leftmargin=*]
    \item \texttt{main.py} - Programa principal
    \item \texttt{utilidades\_loteria.py} - Módulo personalizado
    \item \texttt{requirements.txt} - Dependencias
    \item \texttt{README.txt} - Instrucciones de instalación y uso
    \item Carpeta \texttt{datos\_loteria/} se crea automáticamente
\end{enumerate}

\section{Defensa del Código (60\% de la nota final)}

\subsection{¿Qué es la defensa del código?}

La defensa del código es una reunión virtual donde presentarás y explicarás tu trabajo. Esta defensa es \textbf{obligatoria} y representa el \textbf{60\% de la nota final} de esta tarea.

\subsection{¿Cómo funciona?}

Durante la defensa deberás:

\begin{enumerate}[leftmargin=*]
    \item \textbf{Compartir tu pantalla} mostrando tu código y estructura de archivos.
    
    \item \textbf{Demostrar instalación de dependencias:}
    \begin{itemize}
        \item Mostrar archivo \texttt{requirements.txt}
        \item Ejecutar \texttt{pip install -r requirements.txt}
        \item Verificar que colorama se instala correctamente
    \end{itemize}
    
    \item \textbf{Ejecutar el programa completo:}
    \begin{itemize}
        \item Comprar billetes
        \item Realizar sorteo
        \item Verificar premios
        \item Mostrar estadísticas
        \item Ver información del sistema
    \end{itemize}
    
    \item \textbf{Explicar el uso de cada módulo:}
    \begin{itemize}
        \item Por qué usaste \texttt{random} y qué funciones
        \item Cómo usaste \texttt{statistics} para cálculos
        \item Para qué usaste \texttt{datetime}
        \item Qué haces con \texttt{os} y \texttt{sys}
        \item Cómo creaste y usas tu módulo personalizado
    \end{itemize}
    
    \item \textbf{Responder preguntas como:}
    \begin{itemize}
        \item ``¿Cuál es la diferencia entre \texttt{import random} y \texttt{from random import randint}?''
        \item ``¿Qué hace \texttt{random.seed()}? ¿Para qué serviría?''
        \item ``¿Cómo instalaste colorama?''
        \item ``¿Qué pasa si colorama no está instalado?''
        \item ``¿Por qué crear un módulo propio?''
        \item ``¿Cómo funciona la importación de módulos en Python?''
        \item ``¿Qué es \texttt{requirements.txt}?''
    \end{itemize}
    
    \item \textbf{Demostrar comprensión de:}
    \begin{itemize}
        \item Diferencia entre módulos de biblioteca estándar y externos
        \item Cómo instalar paquetes con pip
        \item Cómo crear y estructurar módulos propios
        \item Buenas prácticas de importación
    \end{itemize}
\end{enumerate}

\subsection{Preguntas Típicas de la Defensa}

Prepárate para responder:

\begin{itemize}[leftmargin=*]
    \item ¿Cuál es la diferencia entre \texttt{import random} y \texttt{from random import *}?
    \item ¿Qué es pip y para qué sirve?
    \item ¿Cómo verificas si un módulo está instalado?
    \item ¿Qué hace \texttt{if \_\_name\_\_ == "\_\_main\_\_"}?
    \item ¿Por qué usar \texttt{try-except} al importar colorama?
    \item ¿Qué es un docstring y por qué usarlo?
    \item ¿Cómo se organiza un proyecto con múltiples archivos?
    \item Dame 3 funciones de \texttt{random} y para qué sirven
    \item ¿Qué diferencia hay entre \texttt{mean} y \texttt{median}?
\end{itemize}

\subsection{Consejos para una Buena Defensa}

\begin{itemize}[leftmargin=*]
    \item \textbf{Conoce cada módulo:} Lee la documentación de los que usaste
    \item \textbf{Instala antes:} Verifica que pip funciona y colorama se instala
    \item \textbf{Explica tu módulo:} Entiende cada función que creaste
    \item \textbf{Ten ejemplos:} Muestra diferentes funcionalidades en acción
    \item \textbf{Entiende imports:} Sé claro sobre las diferencias de importación
\end{itemize}

\subsection{Distribución de la Nota}

\begin{itemize}
    \item \textbf{Código funcional (40\%):} El programa usa correctamente múltiples librerías
    \item \textbf{Defensa del código (60\%):} Tu capacidad de explicar módulos y librerías
\end{itemize}

\textbf{Importante:} Si no asistes a la defensa o no puedes explicar las librerías que usaste, la nota máxima será 40/100.

\section{Desafíos Opcionales (Puntos Extra)}

\begin{enumerate}[leftmargin=*]
    \item \textbf{Módulo JSON para persistencia (+6 pts):}
    \begin{itemize}
        \item Usar \texttt{json} para guardar/cargar datos
        \item Guardar historial de sorteos en formato JSON
        \item Cargar datos anteriores al iniciar
    \end{itemize}
    
    \item \textbf{Gráficos con matplotlib (+8 pts):}
    \begin{itemize}
        \item Instalar: \texttt{pip install matplotlib}
        \item Generar gráfico de barras de números más jugados
        \item Gráfico de pastel de distribución de premios
        \item Histograma de números ganadores históricos
    \end{itemize}
    
    \item \textbf{Módulo collections (+5 pts):}
    \begin{itemize}
        \item Usar \texttt{Counter} para contar frecuencias
        \item Usar \texttt{defaultdict} para organizar datos
        \item Mostrar top 10 números más jugados
    \end{itemize}
    
    \item \textbf{Módulo argparse (+5 pts):}
    \begin{itemize}
        \item Agregar argumentos de línea de comandos
        \item \texttt{python main.py --sorteo} para hacer sorteo directo
        \item \texttt{python main.py --stats} para ver estadísticas
    \end{itemize}
    
    \item \textbf{Módulo pickle para datos binarios (+4 pts):}
    \begin{itemize}
        \item Guardar objetos complejos con pickle
        \item Serializar y deserializar datos
        \item Comparar con JSON
    \end{itemize}
    
    \item \textbf{Crear paquete completo (+7 pts):}
    \begin{itemize}
        \item Organizar en múltiples módulos
        \item Crear \_\_init\_\_.py para paquete
        \item Estructura: loteria/core/, loteria/utils/, loteria/data/
    \end{itemize}
\end{enumerate}

\textit{Nota: Los puntos extra solo se otorgan si la implementación es correcta Y puedes explicarla en la defensa.}

\newpage

\section{Rúbrica de Evaluación}

\subsection{Distribución General}

\begin{itemize}
    \item \textbf{Código Funcional:} 40\% (40 puntos)
    \item \textbf{Defensa del Código:} 60\% (60 puntos)
    \item \textbf{Total:} 100 puntos
    \item \textbf{Puntos Extra:} Hasta +35 puntos adicionales
\end{itemize}

\subsection{Parte 1: Código Funcional (40 puntos)}

\subsubsection{1.1 Uso de Módulos Estándar (20 puntos)}

\begin{itemize}[leftmargin=*]
    \item \textbf{Módulo random (5 pts):}
    \begin{itemize}
        \item 5 pts: Usa 3+ funciones diferentes apropiadamente
        \item 4 pts: Usa 2 funciones correctamente
        \item 2 pts: Usa solo 1 función
        \item 0 pts: No usa random o incorrectamente
    \end{itemize}
    
    \item \textbf{Módulo statistics (5 pts):}
    \begin{itemize}
        \item 5 pts: Usa 3+ funciones para análisis de datos
        \item 4 pts: Usa 2 funciones correctamente
        \item 2 pts: Uso mínimo
        \item 0 pts: No usa statistics
    \end{itemize}
    
    \item \textbf{Módulo datetime (4 pts):}
    \begin{itemize}
        \item 4 pts: Usa para fechas, horas, formateo y cálculos
        \item 3 pts: Uso correcto pero limitado
        \item 1 pt: Uso mínimo
        \item 0 pts: No usa datetime
    \end{itemize}
    
    \item \textbf{Módulos os y sys (3 pts):}
    \begin{itemize}
        \item 3 pts: Usa ambos para interacción con sistema
        \item 2 pts: Usa uno correctamente
        \item 1 pt: Uso mínimo
        \item 0 pts: No usa ninguno
    \end{itemize}
    
    \item \textbf{Importaciones correctas (3 pts):}
    \begin{itemize}
        \item 3 pts: Todas las importaciones bien organizadas y apropiadas
        \item 2 pts: Importaciones funcionales con errores menores
        \item 1 pt: Importaciones desorganizadas
        \item 0 pts: Importaciones incorrectas
    \end{itemize}
\end{itemize}

\subsubsection{1.2 Módulo Personalizado (10 puntos)}

\begin{itemize}[leftmargin=*]
    \item \textbf{Creación del módulo (4 pts):}
    \begin{itemize}
        \item 4 pts: Archivo separado con funciones bien definidas
        \item 3 pts: Módulo funcional con errores menores
        \item 1 pt: Intento de módulo pero incompleto
        \item 0 pts: No crea módulo propio
    \end{itemize}
    
    \item \textbf{Funciones en el módulo (3 pts):}
    \begin{itemize}
        \item 3 pts: 4+ funciones útiles y bien documentadas
        \item 2 pts: 2-3 funciones
        \item 1 pt: Solo 1 función
        \item 0 pts: Sin funciones útiles
    \end{itemize}
    
    \item \textbf{Uso del módulo (3 pts):}
    \begin{itemize}
        \item 3 pts: Importa y usa correctamente en programa principal
        \item 2 pts: Importa pero uso limitado
        \item 1 pt: Importación con errores
        \item 0 pts: No usa su propio módulo
    \end{itemize}
\end{itemize}

\subsubsection{1.3 Funcionalidad del Programa (10 puntos)}

\begin{itemize}[leftmargin=*]
    \item \textbf{Compra de billetes (3 pts):}
    \begin{itemize}
        \item 3 pts: Sistema completo con validaciones
        \item 2 pts: Funcional con errores menores
        \item 1 pt: Implementación básica
        \item 0 pts: No funciona
    \end{itemize}
    
    \item \textbf{Sistema de sorteo (3 pts):}
    \begin{itemize}
        \item 3 pts: Sorteo completo con todos los premios
        \item 2 pts: Sorteo funcional parcialmente
        \item 1 pt: Sorteo básico
        \item 0 pts: No implementa sorteo
    \end{itemize}
    
    \item \textbf{Estadísticas (2 pts):}
    \begin{itemize}
        \item 2 pts: Estadísticas completas con statistics
        \item 1 pt: Estadísticas básicas
        \item 0 pts: No calcula estadísticas
    \end{itemize}
    
    \item \textbf{Persistencia de datos (2 pts):}
    \begin{itemize}
        \item 2 pts: Usa os para crear carpetas y guardar datos
        \item 1 pt: Guarda algunos datos
        \item 0 pts: No guarda datos
    \end{itemize}
\end{itemize}

\subsection{Parte 2: Defensa del Código (60 puntos)}

\subsubsection{2.1 Comprensión de Módulos (30 puntos)}

\begin{itemize}[leftmargin=*]
    \item \textbf{Explicación de importaciones (8 pts):}
    \begin{itemize}
        \item 8 pts: Explica perfectamente import, from...import, diferencias
        \item 6 pts: Explica correctamente con pequeñas dudas
        \item 3 pts: Explicación básica
        \item 0 pts: No entiende importaciones
    \end{itemize}
    
    \item \textbf{Explicación de módulos usados (12 pts):}
    \begin{itemize}
        \item 12 pts: Explica claramente qué hace cada módulo y por qué lo usó
        \item 9 pts: Explica la mayoría correctamente
        \item 5 pts: Explicación superficial
        \item 0 pts: No entiende los módulos que usó
    \end{itemize}
    
    \item \textbf{Instalación de paquetes externos (6 pts):}
    \begin{itemize}
        \item 6 pts: Explica pip, requirements.txt, instalación correctamente
        \item 4 pts: Conoce el proceso con pequeñas dudas
        \item 2 pts: Conocimiento básico
        \item 0 pts: No sabe instalar paquetes
    \end{itemize}
    
    \item \textbf{Creación de módulos propios (4 pts):}
    \begin{itemize}
        \item 4 pts: Explica perfectamente cómo y por qué creó su módulo
        \item 3 pts: Explicación correcta
        \item 1 pt: Explicación vaga
        \item 0 pts: No entiende su propio módulo
    \end{itemize}
\end{itemize}

\subsubsection{2.2 Demostración y Razonamiento (20 puntos)}

\begin{itemize}[leftmargin=*]
    \item \textbf{Demostración de funcionalidades (8 pts):}
    \begin{itemize}
        \item 8 pts: Demuestra todas las funcionalidades exitosamente
        \item 6 pts: Demuestra la mayoría
        \item 3 pts: Demostración parcial
        \item 0 pts: No puede demostrar funcionamiento
    \end{itemize}
    
    \item \textbf{Instalación de dependencias (4 pts):}
    \begin{itemize}
        \item 4 pts: Demuestra instalación con pip y requirements.txt
        \item 3 pts: Demuestra pero con errores menores
        \item 1 pt: Tiene dificultades
        \item 0 pts: No puede instalar dependencias
    \end{itemize}
    
    \item \textbf{Respuesta a preguntas técnicas (8 pts):}
    \begin{itemize}
        \item 8 pts: Responde correctamente todas las preguntas sobre módulos
        \item 6 pts: Responde la mayoría correctamente
        \item 3 pts: Respuestas parciales
        \item 0 pts: No puede responder preguntas básicas
    \end{itemize}
\end{itemize}

\subsubsection{2.3 Calidad y Organización (10 puntos)}

\begin{itemize}[leftmargin=*]
    \item \textbf{Estructura del proyecto (4 pts):}
    \begin{itemize}
        \item 4 pts: Archivos bien organizados, estructura clara
        \item 3 pts: Organización aceptable
        \item 1 pt: Desorganizado
        \item 0 pts: Sin estructura
    \end{itemize}
    
    \item \textbf{Documentación (3 pts):}
    \begin{itemize}
        \item 3 pts: README, docstrings, comentarios claros
        \item 2 pts: Documentación básica
        \item 1 pt: Documentación mínima
        \item 0 pts: Sin documentación
    \end{itemize}
    
    \item \textbf{Requirements.txt (3 pts):}
    \begin{itemize}
        \item 3 pts: Archivo correcto con todas las dependencias
        \item 2 pts: Archivo presente con errores menores
        \item 1 pt: Incompleto
        \item 0 pts: No incluye requirements.txt
    \end{itemize}
\end{itemize}

\subsection{Puntos Extra (Hasta +35 puntos)}

\begin{itemize}[leftmargin=*]
    \item \textbf{Módulo JSON (+6 pts):} Persistencia de datos en formato JSON
    \item \textbf{Gráficos con matplotlib (+8 pts):} Visualizaciones de estadísticas
    \item \textbf{Módulo collections (+5 pts):} Counter, defaultdict para análisis
    \item \textbf{Módulo argparse (+5 pts):} Argumentos de línea de comandos
    \item \textbf{Módulo pickle (+4 pts):} Serialización binaria de datos
    \item \textbf{Crear paquete completo (+7 pts):} Estructura con \_\_init\_\_.py
\end{itemize}

\textbf{Nota:} Los puntos extra solo se otorgan si la implementación es correcta Y el estudiante puede explicarla completamente durante la defensa.

\subsection{Penalizaciones}

\begin{itemize}[leftmargin=*]
    \item \textbf{-100\% de la defensa:} No asiste sin justificación (nota máxima = 40/100)
    \item \textbf{-50\% de la defensa:} Evidencia de no entender los módulos usados
    \item \textbf{-5 pts:} Entrega tardía de 1 día
    \item \textbf{-10 pts:} Entrega tardía de 2-3 días
    \item \textbf{-20 pts:} Entrega tardía de 4+ días
    \item \textbf{-5 pts:} No incluye módulo personalizado
    \item \textbf{-3 pts:} No incluye requirements.txt
    \item \textbf{-3 pts:} Usa \texttt{from modulo import *} (mala práctica)
    \item \textbf{-5 pts:} No usa al menos 3 módulos diferentes
\end{itemize}

\subsection{Nota Mínima Aprobatoria}

Para aprobar esta tarea necesitas obtener al menos \textbf{60 puntos de 100}.

\subsection{Escala de Calificación}

\begin{itemize}
    \item \textbf{90-100+ pts:} Excelente - Dominio completo de módulos y librerías
    \item \textbf{80-89 pts:} Muy Bueno - Comprensión sólida del uso de librerías
    \item \textbf{70-79 pts:} Bueno - Comprensión aceptable, necesita práctica
    \item \textbf{60-69 pts:} Suficiente - Comprensión básica, necesita refuerzo
    \item \textbf{0-59 pts:} Insuficiente - No aprobado
\end{itemize}

\section{Recursos de Apoyo}

\begin{itemize}[leftmargin=*]
    \item \textbf{Documentación oficial:} \url{https://docs.python.org/es/3/tutorial/modules.html}
    \item \textbf{PyPI (paquetes):} \url{https://pypi.org/}
    \item \textbf{Módulo random:} \url{https://docs.python.org/es/3/library/random.html}
    \item \textbf{Módulo statistics:} \url{https://docs.python.org/es/3/library/statistics.html}
    \item \textbf{Módulo datetime:} \url{https://docs.python.org/es/3/library/datetime.html}
    \item \textbf{Guía de pip:} \url{https://pip.pypa.io/en/stable/user_guide/}
\end{itemize}

\section{Preguntas Frecuentes}

\textbf{P: ¿Cuál es la diferencia entre \texttt{import random} y \texttt{from random import randint}?}\\
R: Con \texttt{import random} usas \texttt{random.randint()}. Con \texttt{from random import randint} usas solo \texttt{randint()}.

\textbf{P: ¿Qué es pip?}\\
R: Es el gestor de paquetes de Python. Permite instalar librerías externas.

\textbf{P: ¿Cómo instalo colorama?}\\
R: Ejecuta en la terminal: \texttt{pip install colorama}

\textbf{P: ¿Qué pasa si no tengo pip?}\\
R: Pip viene incluido con Python 3.4+. Si no lo tienes, instala Python de nuevo.

\textbf{P: ¿Puedo usar otros módulos no mencionados?}\\
R: Sí! Mientras cumplas con los requeridos, puedes usar otros adicionales.

\textbf{P: ¿Cómo creo un módulo propio?}\\
R: Simplemente crea un archivo .py con funciones, luego importalo con \texttt{import nombre\_archivo}

\textbf{P: ¿Qué es \_\_name\_\_ == "\_\_main\_\_"?}\\
R: Permite ejecutar código solo cuando el archivo se ejecuta directamente, no cuando se importa.

\vspace{1cm}

\begin{center}
\textbf{¡Éxito con tu simulador de lotería!}\\
\textit{Los módulos y librerías son el superpoder de Python.\\
Aprende a usarlos y tendrás herramientas para cualquier problema.}
\end{center}

\end{document}
